\begin{figure}[!hbtp]
    \centering
    \resizebox{\columnwidth}{!}{%
    \begin{tikzpicture}[
        main_block/.style={
            rectangle, rounded corners=3pt, inner sep=5pt, align=center,
            font=\small\bfseries, draw=black!70, fill=blue!10, line width=0.8pt
        },
        title_block/.style={
            rectangle, rounded corners=3pt, inner sep=3.5pt, align=center,
            font=\footnotesize\bfseries, draw=black!70, fill=blue!10, line width=0.8pt
        },
        arrow/.style = {thick, ->, >={Stealth[length=2.5mm]}},
        box_delimitador/.style = {
            rectangle, rounded corners=8pt, draw=black!70, fill=black!2, line width=0.8pt
        },
        model/.style = {
            rectangle, rounded corners=3pt, align=center, draw=black!70, line width=0.7pt,
            font=\footnotesize, minimum width=2.5cm, minimum height=0.75cm, inner sep=3pt
        },
        legend_sq/.style = {
            rectangle, draw=black!70, line width=0.5pt, minimum size=0.28cm, inner sep=0pt
        },
        cnn/.style = {fill=blue!15},
        trf/.style = {fill=green!15},
        hyb/.style = {fill=orange!22},
        rec/.style = {fill=blue!15},
        base/.style = {fill=black!10}
    ]

        %%%%%%%%%%%%%%%% (a) VIA TOPOLÓGICA %%%%%%%%%%%%%%%%
        \node (ds_a)  [main_block, minimum width=1.9cm] at (0,0) {\textit{Datasets}};
        \node (ssn_a) [main_block, minimum width=2.2cm] at (3.1,0) {SSN};

        \node (grp_a) [box_delimitador, minimum width=9.0cm, minimum height=3.6cm] at (9.6,0) {};
        \node (grp_a_t) [title_block] at (grp_a.north) {6 modelos testados individualmente};
        \node [model, rec,  minimum width=2.5cm] at (6.9,0.55)  {LSTM};
        \node [model, rec,  minimum width=2.5cm] at (9.6,0.55)  {BiLSTM};
        \node [model, rec,  minimum width=2.5cm] at (12.3,0.55) {GRU};
        \node [model, rec,  minimum width=2.5cm] at (6.9,-0.45) {BiGRU};
        \node [model, rec,  minimum width=2.5cm] at (9.6,-0.45) {SAC-LSTM};
        \node [model, base, minimum width=2.5cm] at (12.3,-0.45) {MLP};
        \node [legend_sq, rec]  at (6.2,-1.35) {};
        \node [anchor=west, font=\scriptsize] at (6.4,-1.35) {Recorrente};
        \node [legend_sq, base] at (8.9,-1.35) {};
        \node [anchor=west, font=\scriptsize] at (9.1,-1.35) {\textit{Baseline} não recorrente};

        \node (head_a) [main_block, minimum width=2.4cm, align=center] at (16.2,0) {Cabeça\\MLP};
        \node (met_a)  [main_block, minimum width=1.9cm] at (19.3,0) {Métricas};

        \draw [arrow] (ds_a.east)  -- (ssn_a.west);
        \draw [arrow] (ssn_a.east) -- (grp_a.west);
        \node [below=2pt of ssn_a, font=\scriptsize] {$15\times18$};
        \draw [arrow] (grp_a.east) -- (head_a.west);
        \draw [arrow] (head_a.east) -- (met_a.west);

        \node [font=\small] at (9.65,-2.35) {(a) Via topológica};

        %%%%%%%%%%%%%%%% (b) VIA VISUAL %%%%%%%%%%%%%%%%
        \begin{scope}[yshift=-6.4cm]
        \node (ds_b)  [main_block, minimum width=1.9cm] at (0,0) {\textit{Datasets}};
        \node (aug_b) [main_block, minimum width=2.2cm, align=center] at (3.1,0) {\footnotesize Aumento de\\[-1pt]\footnotesize dados};

        \node (grp_b) [box_delimitador, minimum width=9.0cm, minimum height=3.6cm] at (9.6,0) {};
        \node (grp_b_t) [title_block] at (grp_b.north) {6 modelos testados individualmente};
        \node [model, cnn] at (6.9,0.55)  {ResNet-50};
        \node [model, cnn] at (6.9,-0.45) {EfficientNet-B0};
        \node [model, trf] at (9.6,0.55)  {ViT-Tiny};
        \node [model, trf] at (9.6,-0.45) {DeiT-Tiny};
        \node [model, hyb] at (12.3,0.55) {MobileViT-S};
        \node [model, hyb] at (12.3,-0.45){CoAtNet-0};
        \node [legend_sq, cnn] at (5.9,-1.35) {};
        \node [anchor=west, font=\scriptsize] at (6.1,-1.35) {Convolucional};
        \node [legend_sq, trf] at (8.4,-1.35) {};
        \node [anchor=west, font=\scriptsize] at (8.6,-1.35) {\textit{Transformer}};
        \node [legend_sq, hyb] at (11.0,-1.35) {};
        \node [anchor=west, font=\scriptsize] at (11.2,-1.35) {Híbrida};

        \node (head_b) [main_block, minimum width=2.4cm, align=center] at (16.2,0) {Cabeça\\MLP};
        \node (met_b)  [main_block, minimum width=1.9cm] at (19.3,0) {Métricas};

        \draw [arrow] (ds_b.east)  -- (aug_b.west);
        \draw [arrow] (aug_b.east) -- (grp_b.west);
        \node [below=2pt of aug_b, font=\scriptsize] {$256\times256$};
        \draw [arrow] (grp_b.east) -- (head_b.west);
        \draw [arrow] (head_b.east) -- (met_b.west);

        \node [font=\small] at (9.65,-2.35) {(b) Via visual};
        \end{scope}

    \end{tikzpicture}%
    }
    \caption{Ablação das vias isoladas. (a) Via topológica: a sequência de características da SSN é processada por seis codificadores (cinco recorrentes e um MLP como \textit{baseline}), cada um testado de forma independente. (b) Via visual: as imagens, após aumento de dados, são processadas por seis modelos (convolucionais, \textit{transformers} e híbridos), também testados individualmente. Cada via possui sua própria cabeça MLP, da qual são extraídas as métricas.}
    \label{fig:ablacao-vias}
\end{figure}
